\chapter{System Architecture}
\label{ch:arch}

\section{Overview}

HALQA is a three-tier system: a React single-page PWA, a NestJS REST API, and MongoDB.
Its distinguishing architectural property is not the layering but the \emph{isolation of
the model provider}: exactly one file in the backend knows a model name or an API key,
and a unit test enforces that no other file imports it. Everything else --- routing,
quota, caching, degradation, grounding --- is ordinary, testable TypeScript that happens
to orchestrate model calls.

\begin{figure}[htbp]
\centering
\begin{tikzpicture}[
  font=\sffamily\footnotesize, x=1cm, y=1cm,
  lay/.style={draw=ink0,fill=white,line width=0.5pt,rounded corners=2pt,
              text width=11.4cm,minimum height=1.4cm,align=left,
              inner sep=6pt},
  fade/.style={draw=ink2,fill=fillA,line width=0.4pt,rounded corners=2pt,
               text width=11.4cm,minimum height=1.1cm,align=left,
               inner sep=5pt},
  tag/.style={font=\sffamily\bfseries\scriptsize,color=ink3},
  ar/.style={-{Latex[length=1.5mm,width=1.3mm]},draw=ink0,line width=0.6pt},
]
\node[lay] (t1) at (0,0) {\textbf{React 19 PWA} \hfill{\color{ink3}15 routes, 4 registers}
  \\[2pt]{\footnotesize precached shell \enspace\textbullet\enspace service worker
  \enspace\textbullet\enspace session store \enspace\textbullet\enspace offline pack}};
\draw[ar] (t1.south) -- ++(0,-0.42) node[right=3pt,font=\sffamily\tiny,color=ink2]{};
\node[fade] (t2) at (0,-2.05) {\textbf{NestJS 11 API} \hfill{\color{ink3}20 controllers, 68 routes}
  \\[2pt]{\footnotesize guard chain (token $\to$ role $\to$ rate) \enspace\textbullet\enspace
  Zod-validated boot \enspace\textbullet\enspace uniform error contract}};
\draw[ar] (t2.south) -- ++(0,-0.42) node[right=3pt]{};
\node[lay] (t3) at (0,-4.10) {\textbf{AI platform} \hfill{\color{ink3}one provider adapter}
  \\[2pt]{\footnotesize AI router $\to$ semantic cache $\to$ quota governor $\to$ FIFO queue
  $\to$ adapter (C / L) \enspace\textbullet\enspace TF-IDF retrieval, 39 chunks
  \enspace\textbullet\enspace degradation green--black}};
\draw[ar] (t3.south) -- ++(0,-0.42) node[right=3pt]{};
\node[fade] (t4) at (0,-6.15) {\textbf{MongoDB 8} \hfill{\color{ink3}18 collections, 22 indexes}
  \\[2pt]{\footnotesize 90-day TTL on telemetry \enspace\textbullet\enspace
  atomic quota reservation \enspace\textbullet\enspace two idempotent migrations}};
% side annotations
\node[font=\sffamily\tiny,color=ink3,anchor=east,align=right] at (-6.45,0) {TIER 1\\CLIENT};
\node[font=\sffamily\tiny,color=ink3,anchor=east,align=right] at (-6.45,-2.05) {TIER 2\\API};
\node[font=\sffamily\tiny,color=ink3,anchor=east,align=right] at (-6.45,-4.10) {AI\\PLATFORM};
\node[font=\sffamily\tiny,color=ink3,anchor=east,align=right] at (-6.45,-6.15) {TIER 3\\DATA};
\end{tikzpicture}
\caption{The four-layer stack, read top to bottom. Each layer states what it owns and the
one property it is responsible for keeping true. Model calls never skip a layer: they enter
at the AI platform and are metered, cached and grounded before the adapter is reached.}
\label{fig:topology}
\end{figure}

\section{Layered decomposition}

The backend is organised so that most of its logic is testable without a database, a
network or a Nest container. Two layers carry that responsibility.

\paragraph{\code{domain/} --- pure logic.} Fifteen modules with no framework imports:
PII screening, credit mathematics with its logarithmic cap, S-B-U health formulas, cache
key construction, the Arabizi normalizer, the bandit update, hint templates, token
budgeting, confusion signals, ID formats and context modes. Because these functions are
pure, they are tested directly and exhaustively --- 145 unit assertions across 24 spec
files --- and a refactor of the surrounding services cannot silently change them.

\paragraph{\code{modules/} --- use cases.} Twenty-five feature modules, each owning a
controller, a service and its schema access. A module may depend on \code{domain/}, on
the event bus, and on other feature services, but never on the provider adapter. The
adapter-isolation test enforces that dependency rule mechanically:

\begin{lstlisting}[language=Java,caption={The architectural gate that keeps the provider from leaking into product code (\code{test/unit/adapter-isolation.spec.ts}).}]
it('only src/ai files import the provider adapter', () => {
  const offenders: string[] = [];
  for (const f of tsFiles(SRC)) {
    if (f.includes(`${path.sep}ai${path.sep}`)) continue;
    const src = fs.readFileSync(f, 'utf8');
    if (/(?:import|require)\s*\(?[^'"]*['"][^'"]*ai\/adapter\/provider\.adapter/.test(src)) {
      offenders.push(path.relative(SRC, f));
    }
  }
  expect(offenders).toEqual([]);
});
\end{lstlisting}

This is the difference between a claim and a constraint. The report can say ``the
provider is isolated in one adapter'' because a failing test says so on every build.

\section{The two model categories}
\label{sec:routing}

The challenge expects participation in both model categories. In an environment without
GPU acceleration and with shared access, there are two symmetric failure modes: sending
everything to the larger model, which exhausts the scarce quota and is slower for no
benefit; or capping everything at the compact model, which produces thin synthesis where
synthesis is genuinely required. HALQA makes the split explicit, encodes it in
configuration, and measures it.

\begin{table}[htbp]
\centering
\caption{Routing policy. The compact tier is the default; the larger tier must earn the call.}
\label{tab:routing}
\footnotesize
\begin{tabularx}{\textwidth}{@{}p{3.0cm}X X@{}}
\toprule
& \textbf{Compact (C) --- default} & \textbf{Larger (L) --- escalation only} \\
\midrule
Tasks &
  Group Pulse routing (intent, register, safety, difficulty); Quick Take drafts; hint
  ladder L1--L3; poll and quiz items; language rewriting; DNA interpretation; Whisper
  compare; mentor briefs &
  $\ge 3$ retrieved chunks in conflict (multi-source synthesis); polished capsule or
  revision prose; low C confidence; explicit Whisper or Auditor escalation; deep-dive
  question \\
Cost profile &
  One schema-constrained call per turn window, shared by credits, confusion, intent and
  register &
  Seconds of latency; counted against a separate, smaller quota \\
Guardrail &
  Falls back to template (T0) where a template exists &
  Falls back to C, then to T0, then to the offline pack \\
\bottomrule
\end{tabularx}
\end{table}

The policy is not prose in a design document. It lives in \code{config/router.yaml} as
thresholds, and its outcome is measured from the persisted call ledger:

\begin{lstlisting}[caption={\code{config/router.yaml} --- the thresholds that the router reads at boot.}]
confidence_threshold: 0.55     # below this, a C verdict is not trusted
synthesis_chunk_conflict: 3    # >= 3 conflicting chunks forces synthesis
l_share_target: 0.25           # ceiling on larger-tier share of all calls
skeleton_ms: 200               # zero-model-call first paint budget
c_complete_ms: 4000            # compact completion budget
l_complete_ms: 12000           # larger completion budget
\end{lstlisting}

The escalation gate itself is deliberately conservative, and is expressed as an ordered
list of conditions rather than a score, so that a reviewer can predict its behaviour:

\begin{lstlisting}[language=Java,caption={Escalation conditions in the router, in evaluation order.}]
//  1. >= 3 chunks with conflicting claims need synthesis
//  2. the compact classifier is not confident about the intent
//  3. the learner explicitly asked for a deep dive
//  4. an Auditor or Whisper verdict requested escalation
//  Anything else: compact tier, or a T0 template with no call at all.
\end{lstlisting}

\section{Request lifecycle: the Circle-First Ask}
\label{sec:cfas}

The Circle-First Ask is the system's spine, and every other feature reuses one of its
stages. The sequence is deliberately asymmetric: the AI does its expensive work
\emph{while the humans are talking}, so the peer window costs the learner nothing in
latency.

\begin{figure}[htbp]
\centering
\begin{tikzpicture}[
  font=\sffamily\scriptsize,
  act/.style={draw=ink2,fill=white,rounded corners=1.5pt,minimum height=9.5mm,
              minimum width=29mm,align=center,line width=0.4pt},
  wact/.style={draw=ink0,fill=fillB,rounded corners=1.5pt,minimum height=9.5mm,
              minimum width=29mm,align=center,line width=0.5pt},
  ai/.style={draw=ink0,fill=fillC,text=white,rounded corners=1.5pt,minimum height=9.5mm,
             minimum width=30mm,align=center,line width=0.5pt},
  ar/.style={-{Latex[length=1.3mm,width=1.1mm]},draw=ink0,line width=0.5pt},
  dr/.style={-{Latex[length=1.3mm,width=1.1mm]},draw=ink2,line width=0.4pt,dashed},
  lane/.style={font=\sffamily\footnotesize\bfseries,color=ink1},
  tm/.style={font=\sffamily\tiny,color=ink2},
]
% lanes
\node[lane,anchor=west] at (-0.4,0) {Learner};
\node[lane,anchor=west] at (-0.4,-1.75) {Peers};
\node[lane,anchor=west] at (-0.4,-3.5) {Nova};
\node[lane,anchor=west] at (-0.4,-5.25) {Server};

% learner
\node[wact] (l1) at (1.6,0) {post question};
\node[act]  (l2) at (5.4,0) {watch window};
\node[act]  (l3) at (9.0,0) {read verified};
\node[act]  (l4) at (12.2,0) {vote / vote-fix};

% peers
\node[wact] (p1) at (5.4,-1.75) {answer in window};
\node[act]  (p2) at (9.0,-1.75) {see verdict};

% nova
\node[ai] (n1) at (5.4,-3.5) {retrieve $\to$ cache $\to$ C draft\\(hidden)};
\node[ai] (n2) at (9.0,-3.5) {Auditor verify\\(1$\times$C, poll-first)};
\node[ai] (n3) at (12.2,-3.5) {release with cites};

% server
\node[act] (s1) at (1.6,-5.25) {open window\\ \code{OPEN}};
\node[act] (s2) at (7.2,-5.25) {expiry worker\\ settles};
\node[act] (s3) at (12.2,-5.25) {\code{RESOLVED}\\ + credits, + S-B-U};

% timeline
\draw[ink4,line width=0.4pt] (0.2,-6.4) -- (13.3,-6.4);
\foreach \x/\t in {1.6/0\,s, 5.4/1\,s, 7.2/45\,s, 9.0/46\,s, 12.2/48\,s}{
  \draw[ink4] (\x,-6.4) -- (\x,-6.6);
  \node[tm] at (\x,-6.9) {\t};
}
\draw[decorate,decoration={brace,amplitude=4pt,mirror},ink3]
  (7.1,-2.55) -- (3.7,-2.55) node[midway,above=5pt,font=\sffamily\tiny,color=ink2]
  {peer window 45--60\,s, AI work overlapped};

\draw[ar] (l1) -- (l2);
\draw[ar] (l2) -- (l3);
\draw[ar] (l3) -- (l4);
\draw[ar] (l1) -- (s1);
\draw[ar] (l2) -- (p1);
\draw[dr] (p1) -- (s2);
\draw[ar] (s2) -- (n2);
\draw[dr] (n1) -- (n2);
\draw[ar] (n2) -- (n3);
\draw[ar] (n3) -- (l3);
\draw[dr] (n2) -- (p2);
\draw[dr] (n3) -- (s3);
\draw[dr] (l4) -- (s3);
\end{tikzpicture}
\caption{Circle-First Ask. The peer window is the only interval in which humans speak; the
model work is overlapped with it, so the AI answer is ready at the moment the window
closes. The withholding rule is enforced by question state, not by prompt instruction.}
\label{fig:cfas}
\end{figure}

\paragraph{Question state machine.} The withholding rule is a property of persisted state,
which is what makes it trustworthy. A question is \code{OPEN} while peers may answer;
\code{NOVA\_THINKING} after settlement begins, when the Quick Take is already written but
not yet exposed; and \code{RESOLVED} when the verified answer, its citations and any poll
are readable. The private and public views are separate routes, so a client cannot obtain
the hidden answer by reading a field it was not given: the public state route strips it,
and the private route is asker-only.

\paragraph{Pipeline stages.} Progress is reported as an ordered list of stages
(\code{retrieving}, \code{checking\_similar}, \code{writing}, \code{verifying}), each with
a start and completion timestamp. Responses arrive whole rather than streamed, so the
interface communicates staged progress instead of a partially-rendered answer
(\S\ref{sec:latency-ux}).

\section{Why the window is server-authoritative}

A tempting shortcut is to run the countdown in the browser. We rejected it for two
reasons visible in production. First, a client clock is untrusted and a backgrounded tab
throttles timers, so a window could outlive its deadline and peers would keep voting on
a settled question. Second, a single process restart in the middle of a window would
leave the question open forever, which is a visible product defect: the learner sees a
window that never closes.

The implementation therefore keeps the deadline on the server (\code{windowEndsAt} as an
indexed field), settles from a timer, and \emph{also} sweeps periodically so that a window
that expired while the process was down is settled on the next tick rather than never.
Clients derive their countdown from \code{windowEndsAt}, not from a server-provided
``now'', which also keeps conditional-GET validators stable:

\begin{lstlisting}[language=Java,caption={Dual settlement path: timer for the normal case, periodic sweep for the crash case.}]
/** Periodic sweep, so a window that expired while this process was down --
 *  or between a timer being set and a deploy -- is still settled. */
const overdue = await this.questions.find({ state: 'OPEN', windowEndsAt: { $lte: new Date() } });
const upcoming = await this.questions.find({ state: 'OPEN', windowEndsAt: { $gt: new Date() } });
for (const q of upcoming) this.schedule(q._id, +q.windowEndsAt - Date.now());
\end{lstlisting}

The same reasoning applies at deploy time: shutdown hooks drain in-flight requests and
clear the timers, so a deploy does not settle windows from a process that is already
dead. Settlement is idempotent, so a double settle cannot double-credit.

\section{Cross-cutting concerns}

\begin{table}[htbp]
\centering
\caption{Where each cross-cutting concern is implemented, and how it is enforced.}
\label{tab:crosscut}
\footnotesize
\begin{tabularx}{\textwidth}{@{}p{2.7cm}X p{4.3cm}@{}}
\toprule
\textbf{Concern} & \textbf{Implementation} & \textbf{Enforcement} \\
\midrule
Authentication & HMAC bearer token with learner ID, role and a token epoch
  & Default-deny guard; epoch bump invalidates all devices on passphrase change \\
Authorization & Role decorator plus circle-membership check
  & \code{authz-matrix} integration test: learner $\times$ mentor $\times$ admin
    $\times$ anonymous $\times$ cross-circle \\
Rate limiting & Named budgets (AI 10/min, writes 30/min, join-code 60/min per IP)
  & In-memory guard, with the edge limiter documented as the production layer \\
Validation & Zod-validated config at boot; \code{ValidationPipe} with
  \code{whitelist} and \code{forbidNonWhitelisted} on every route
  & Rejects unknown fields rather than ignoring them \\
Error contract & Uniform \code{\{error, code, retryable, user\_msg[4]\}}
  & Localized message asserted for all four languages in unit tests \\
Observability & Request-ID middleware, JSON logger, prom-client registry, event bus
  & \code{/metrics} key set asserted by the live smoke \\
Persistence & Mongoose schemas, 22 indexes, a two-migration registry, unit-of-work
  for transactions
  & Migrations are idempotent and indexed \\
Shutdown & \code{enableShutdownHooks}, timers drained, Mongo closed
  & Deploy smoke check \\
\bottomrule
\end{tabularx}
\end{table}

\section{Technology choices}

Each choice below was made against a constraint rather than by preference, and each is
justified by what it enabled.

\paragraph{NestJS 11 with TypeScript in strict mode.} The feature set is large enough that
module boundaries matter: 25 modules that can depend on \code{domain/} but not on the
adapter is a rule a framework can express and a flat Express app cannot. Strict mode with
no implicit \code{any} was decisive for a codebase that passes \code{unknown} model output
around: every parse of a model response is forced through a narrowing function, which is
where several real bugs were found.

\paragraph{MongoDB 8 with Mongoose.} Session state is document-shaped --- a question
carries its stages, its citations and its Quick Take together --- and a session reads as
one document rather than a join of six tables. Mongo also provides the two operational
properties this project depends on: TTL indexes for telemetry retention and atomic
\code{findOneAndUpdate} for quota reservation, which is what makes the governor correct
across concurrent requests.

\paragraph{React 19 with Vite and a PWA plugin.} A shared-device, campus-network context
wants an installable shell and a precached offline pack. Vite's build plus
\code{vite-plugin-pwa} produced both without a bespoke service worker.

\paragraph{TF-IDF retrieval instead of an embedding index.} This is the most debatable
choice, and it was made deliberately. A 39-chunk knowledge base over two topics does not
need embeddings to be retrieved well, and a local index has three properties a remote
embedding service does not: it builds in under a second at boot, it costs nothing per
query, and it keeps the grounded path alive in the \code{RED} and \code{BLACK} degradation
states. The trade-off is recall on paraphrases, which the three-variant query rewrite
partly compensates (\S\ref{sec:rag}).

\paragraph{A single provider adapter.} All model access goes through one file exposing
\code{compact}, \code{larger} and \code{available(tier)}. Two benefits follow: the product
can survive a provider change or a move to self-hosted open-weight models without touching
feature code, and the isolation is testable rather than aspirational.
